% 370_Warum_Ikosaeder_De_ch.tex
% Johann Pascher, 20. September 2026
% Erklärungsdokument: Drei, Fünf und das Nicht-Kommutieren

\providecommand{\BM}{\textbf{[B]}}
\providecommand{\KM}{\textbf{[K]}}
\providecommand{\SM}{\textbf{[S]}}
\providecommand{\EM}{\textbf{[E]}}

\chapter*{Warum das Ikosaeder: Drei, Fünf und das Nicht-Kommutieren}
\addcontentsline{toc}{chapter}{Warum das Ikosaeder}

\begin{abstract}
FFGFT ist von Grund auf eine $\mathbb{Z}_3$-Theorie. Die Leptonmassen enthalten
aber den goldenen Schnitt $\phi$, und $\phi$ ist die Zahl der Fünf. Dieses
Dokument erklärt, warum das Ikosaeder die Brücke ist: seine Drehgruppe $A_5$ ist
die kleinste endliche Gruppe, in der eine Drei und eine Fünf zusammen vorkommen,
ohne miteinander zu vertauschen. Aus diesem Nicht-Vertauschen entsteht eine
Mischung der $\mathbb{Z}_3$-Moden unter der Fünffachdrehung, und die Mischung
hat eine rationale Komponente $2/9$ und zwei $\phi$-haltige Komponenten. Die
rationale ist die Koide-Phase (Dok.~367), die $\phi$-haltigen sind das Skelett
der Massenverhältnisse (Dok.~368). Ein einziges Matrixelement dieser Mischung
trägt die gesamte Koide-Formel. Das Dokument rechnet nichts Neues; es erzählt,
was die Dokumente 285, 293, 367 und 368 zusammen bedeuten.
\end{abstract}

\section*{Zeichenerklärung}
\addcontentsline{toc}{section}{Zeichenerklärung}

\paragraph*{Epistemische Marker:}\mbox{}\\[2pt]
\begin{tabular}{@{}ll@{}}
\textbf{[B]} & algebraisch bewiesen (Prüfskript im zitierten Dokument) \\
\textbf{[K]} & numerisch bestätigt \\
\textbf{[S]} & strukturell beobachtet, Herleitung offen \\
\textbf{[E]} & etablierte externe Mathematik \\
\end{tabular}

\medskip
\begin{tabular}{@{}lp{11cm}@{}}
\toprule
\textbf{Symbol} & \textbf{Bedeutung} \\
\midrule
$\mathbb{Z}_3$, $C_3$ & zyklische Gruppe der Ordnung 3; $C_3$: Permutation $(x,y,z)\mapsto(y,z,x)$ \\
$v_0,v_1,v_2$ & Fourier-Moden, Eigenvektoren von $C_3$ \\
$\omega$ & $e^{2\pi i/3}$ \\
$A_5$ & Ikosaedergruppe, 60 Elemente \\
$R_5$ & Drehung um $72^\circ$ um die Ikosaeder-Ecke $(0,1,\phi)$ \\
$\phi$ & goldener Schnitt $(1+\sqrt5)/2$ \\
$p_0,p_1,p_2$ & Umverteilungsgewichte der Elektron-Mode unter $R_5$ (Dok.~293) \\
$A$ & $\langle v_0\mid R_5\,v_e\rangle$, $|A|=\sqrt2/3$ (Dok.~367) \\
$D_4$ & Gitter auf $T^4$; $|\text{Aut}(D_4)|=1152$ (Dok.~314, 364) \\
GF$(3^k)$ & endlicher Körper der Charakteristik 3; GF$(9)\ni\phi$, GF$(81)\ni\zeta_5$ \\
\bottomrule
\end{tabular}

\section{Die Beobachtung}

Es gibt drei geladene Leptonen. Ihre Massen stehen im Verhältnis
$m_e : m_\mu : m_\tau \approx 1 : 206{,}8 : 3477{,}4$, und dieses Verhältnis
gehorcht seit 1981 einer verblüffend einfachen Formel, der Koide-Formel:
\[
  \sqrt{m_k} = \sqrt{M}\Bigl(1 + \sqrt2\cos\!\Bigl(\theta + \frac{2\pi k}{3}\Bigr)\Bigr),
  \qquad k = 0,1,2.
\]
Die Skala $M$ ist eine Einheit -- sie sagt nur, wie groß die Massen insgesamt
sind. Die Phase $\theta$ ist die eigentliche Information: sie bestimmt die
Verhältnisse. Empirisch ist $\theta = 2/9$ auf $10^{-7}$ genau (Dok.~292).
Warum $2/9$? Warum nicht $1/4$, nicht $0{,}3$? Und warum der Koeffizient
$\sqrt2$?

Die Antwort führt über ein geometrisches Objekt, das mit Leptonen zunächst
nichts zu tun hat: das Ikosaeder. Der Rest dieses Dokuments erklärt, warum.

\section{Wo FFGFT herkommt: die Drei}

FFGFT beginnt mit einer $\mathbb{Z}_3$. Drei Generationen, Trialität, der
Zirkulant als Massenoperator. Der Modenraum ist $\mathbb{C}^3$, seine
natürliche Basis sind die Fourier-Moden
\[
  v_k = \frac{1}{\sqrt3}\,(1,\ \omega^k,\ \omega^{2k}), \qquad k = 0,1,2.
\]
Das sind die Eigenvektoren der zyklischen Permutation $C_3$. Die Mode $v_0$ ist
die symmetrische; $v_1$ und $v_2$ tragen die Phasen $\omega$ und $\omega^2$. Das
Elektron sitzt in $v_1$, das Myon in $v_2$, das Tau in $v_0$ -- das ist die
Zuordnung, die der Koide-Index $k$ trägt.

Was diese Struktur allein \emph{nicht} liefert, ist eine Zahl. Die $\mathbb{Z}_3$
sagt, dass es drei Moden gibt und wie sie sich unter $C_3$ transformieren. Sie
sagt nichts über Massen. Die Koide-Formel braucht zwei Eingaben -- den
Koeffizienten $\sqrt2$ und die Phase $\theta = 2/9$ -- und beide sind in der
$\mathbb{Z}_3$ nicht enthalten.

\section{Wo $\phi$ herkommt: die Fünf}

Der goldene Schnitt ist die Zahl des Fünfecks: $\cos 72^\circ = (\phi-1)/2$, und
die Diagonalen eines regelmäßigen Fünfecks stehen zu seinen Seiten im
Verhältnis $\phi$ \EM. In der $\mathbb{Z}_3$-Struktur gibt es keine Fünf. Wenn
$\phi$ in den Leptonmassen auftaucht -- Dok.~368 zeigt $p_1/p_2 = \phi^8$ exakt
\BM~--, muss eine Fünf von irgendwoher kommen.

\section{Warum das Ikosaeder: wo Drei und Fünf sich treffen}

Das Ikosaeder hat 20 Dreiecksflächen und 12 Ecken; an jeder Ecke treffen sich
fünf Dreiecke. Seine Drehgruppe $A_5$ hat 60 Elemente: die Identität, 15
Drehungen um $180^\circ$ (Kantenmitten), 20 Drehungen um $120^\circ$
(Flächenmitten), 24 Drehungen um $72^\circ$ oder $144^\circ$ (Ecken) \EM. Drei
und Fünf sind beide da.

Das ist nicht selbstverständlich. Die zyklische Gruppe $\mathbb{Z}_{15}$ enthält
ebenfalls Elemente der Ordnung 3 und 5 -- aber dort vertauschen sie. In $A_5$
\emph{vertauschen sie nicht}: eine Drehung um $120^\circ$, gefolgt von einer
um $72^\circ$, ist etwas anderes als die umgekehrte Reihenfolge. $A_5$ ist die
kleinste Gruppe, in der das so ist.

Konkret: in Standardlage sind die Ecken des Ikosaeders die zyklischen
Permutationen von $(0,\pm1,\pm\phi)$. Die 3-fach-Achsen gehen durch
Flächenmitten; eine davon ist $(1,1,1)$, und die Drehung um $120^\circ$ um diese
Achse ist \emph{genau} die Permutation $C_3$ von FFGFT. Die $\mathbb{Z}_3$ von
FFGFT ist ohne Umdeutung eine Untergruppe der Ikosaedergruppe (Dok.~285). Die
5-fach-Achsen gehen durch Ecken; eine Ecke ist $(0,1,\phi)$, und die Drehung
$R_5$ um $72^\circ$ um diese Achse ist die Fünffachdrehung, die Dok.~293
verwendet.

\section{Wo das Ikosaeder wohnt: ausgerollt, nicht kompakt}

Bis hierher klingt es, als wäre das Ikosaeder einfach \emph{die} Struktur von
FFGFT. Das ist es nicht -- und der Korpus ist darüber sehr genau (Dok.~285,
314, 364). FFGFT hat zwei Bilder, die durch eine verlustfreie
Koordinatentransformation verbunden sind, und das Ikosaeder wohnt nur in
einem davon.

\paragraph{Kompakt: der Torus $T^4$ mit dem $D_4$-Gitter.}
Das ist das Original. Der Vier-Torus trägt das $D_4$-Gitter, und dessen
Trialität ist die $\mathbb{Z}_3$ von FFGFT -- als Gitteroperation, als
Zirkulant auf den Modenorbits (Dok.~314). Die Wicklungszahlen der T0-Leiter
leben hier. Das Gitter kennt die Drei. Es kennt \emph{keine} Fünf:
$|\text{Aut}(D_4)| = 1152 = 2^7\cdot3^2$, und $5\nmid1152$. Kein Ikosaeder
wirkt auf einem $D_4$-Gitter (Dok.~364, Satz zur Unverträglichkeit) \BM.

\paragraph{Ausgerollt: $\mathbb{R}^3\times S^1$.}
Der beobachtete Raum entsteht, indem ein Kreis des Torus durch den
Messvorgang zur Zeit wird: $T^4\to\mathbb{R}^3\times S^1$ (Dok.~285). Auf
\emph{diesem} $\mathbb{R}^3$ wirkt $A_5$ als Drehgruppe. Die $C_3\subset A_5$
ist dieselbe Drei wie die Trialität im kompakten Bild -- die Drehung um
$120^\circ$ um $(1,1,1)$. Die Fünf aber gibt es erst hier. Das Ikosaeder ist
das Bild der Torusgeometrie im ausgerollten Raum, nicht der Torus selbst.

\paragraph{Wie die Fünf trotzdem in den Körper kommt.}
Dok.~364 löst die scheinbare Spannung: es kommt darauf an, \emph{als was}
eine Symmetrie eintritt -- als Automorphismus von Gittervektoren oder als
Phase. Die $72^\circ$-Drehung hat das charakteristische Polynom
$(x-1)(x^2-(\phi-1)x+1)$; ihre Spur ist $\phi$. Modulo~3 liegt $\phi$ in
GF$(9)$ und ist dort primitiv ($\phi^4=-1$, Ordnung~8). Die Wurzeln des
quadratischen Faktors liegen nicht in GF$(9)$, sondern in GF$(81)$, und haben
dort Ordnung~5 -- denn $5\mid80=|\text{GF}(81)^*|$ \BM. Im Gitter ist die
Drehung um $72^\circ$ ein Vektorautomorphismus, und $5\nmid1152$. Im Körper
ist dieselbe Drehung nur noch ihr Eigenwert $\zeta_5$, und $5\mid80$. Die
geometrische Unverträglichkeit wird zur arithmetischen Verträglichkeit, sobald
man von Vektoren zu Phasen übergeht. Dok.~364 wörtlich: das ist der
Mechanismus hinter Dok.~293 -- $\theta=2/9$ entsteht aus $R_5$ als Phase,
nicht als Gitteroperation.

\paragraph{Warum das erzwungen ist.}
In Charakteristik~3 ist $x^3-1=(x-1)^3$ \EM: in keinem GF$(3^k)$ gibt es eine
primitive dritte Einheitswurzel. Die Drei \emph{kann} keine Phase sein und
muss Gitter sein. Die Fünf \emph{kann} kein Gitter sein ($5\nmid1152$) und
muss Phase sein. Die Primzahl~3 verteilt die beiden Symmetrien auf zwei
komplementäre Kanäle -- das ist keine Wahl, sondern Zwang. Die Erweiterung
GF$(9)\to$GF$(81)$, die $\zeta_5$ und $\phi$ trägt, ist genau die Verdopplung
$6=3\oplus3'$ aus Dok.~285.

\paragraph{Was das für die Erzählung bedeutet.}
Wenn im Folgenden vom Nicht-Vertauschen der Drei und der Fünf die Rede ist,
ist gemeint: die Drei als Gitteroperation im kompakten Bild, die Fünf als
Phase im ausgerollten. Beide treffen sich nicht in der Geometrie -- dort
sind sie unverträglich --, sondern im Körper GF$(81)$. Das Ikosaeder ist die
anschauliche Form dieser Begegnung.

\section{Was das Nicht-Vertauschen erzeugt}

Weil $R_5$ nicht mit $C_3$ vertauscht, ist $R_5$ in der Modenbasis
$v_0,v_1,v_2$ \emph{nicht diagonal}. Wendet man $R_5$ auf die Elektron-Mode
$v_1$ an, landet das Ergebnis nicht wieder in $v_1$, sondern wird über alle
drei Moden verteilt. Die Betragsquadrate der Projektionen sind die
Umverteilungsgewichte (Dok.~293) \BM:
\[
  p_0 = |\langle v_0|R_5|v_1\rangle|^2 = \frac{2}{9}, \qquad
  p_1 = |\langle v_1|R_5|v_1\rangle|^2 = \frac{2+3\phi}{9}, \qquad
  p_2 = |\langle v_2|R_5|v_1\rangle|^2 = \frac{5-3\phi}{9}.
\]
Man sieht sofort: $p_0$ ist rational, $p_1$ und $p_2$ enthalten $\phi$. Das ist
Struktur, kein Zufall. Die Projektion auf die symmetrische Mode $v_0$ ist gegen
die Fünffachstruktur blind -- sie zählt nur, wie viel von der
$\mathbb{Z}_3$-Mischung ankommt, und das ist $2/9 = 2/3^2$. Die Projektionen auf
$v_1$ und $v_2$ sehen die Fünf und tragen deshalb $\phi$.

Noch klarer wird es an der symmetrischen Mode selbst. Was $R_5$ aus $v_0$ macht:
\[
  \Bigl(|\langle v_0|R_5|v_0\rangle|^2,\ |\langle v_0|R_5|v_1\rangle|^2,\
  |\langle v_0|R_5|v_2\rangle|^2\Bigr) = \Bigl(\frac{5}{9},\ \frac{2}{9},\ \frac{2}{9}\Bigr) \quad \KM
\]
Die vollständige Matrix der Betragsquadrate \KM:
\begin{center}
\begin{tabular}{@{}l|ccc@{}}
\toprule
$|\langle v_j|R_5|v_k\rangle|^2$ & $k=0$ & $k=1$ & $k=2$ \\
\midrule
$j=0$ & $5/9$ & $2/9$ & $2/9$ \\
$j=1$ & $2/9$ & $(2+3\phi)/9$ & $(5-3\phi)/9$ \\
$j=2$ & $2/9$ & $(5-3\phi)/9$ & $(2+3\phi)/9$ \\
\bottomrule
\end{tabular}
\end{center}
Zwei Dinge sieht man sofort. Überall, wo $v_0$ beteiligt ist, steht $2/9$
oder $5/9$ -- rational, ohne $\phi$. Und $\phi$ erscheint nur zwischen $v_1$
und $v_2$, den beiden Moden, die die Fünf sehen. Jede Zeile und jede Spalte
summiert zu~$1$.

Die Amplituden der ersten Zeile sind $(\sqrt5,\sqrt2,\sqrt2)/3$. Die Fünffachdrehung lässt
$\sqrt5/3$ der symmetrischen Mode in sich selbst und gibt $\sqrt2/3$ an jede der
beiden nicht-trivialen Moden ab. Und $5+2+2 = 9 = 3^2$. Das ist die eine Zeile,
in der die Fünf des Ikosaeders und die Drei von FFGFT nebeneinander stehen. Die
Spur von $R_5$ ist $1 + 2\cos72^\circ = \phi$ \EM~-- die Fünf ist im Operator, die
Drei in der Basis, und das Produkt beider ist die Zeile $(5,2,2)/9$.

\section{Dok.~367: warum $\theta = p_0$, und wie die ganze Koide-Formel folgt}

Dok.~293 hatte festgestellt, dass $p_0 = 2/9$ mit dem Koide-Winkel $\theta = 2/9$
übereinstimmt, und die Frage nach dem Grund offen gelassen. Dok.~367 gibt die
Antwort: beide sind derselbe Quotient $2/3^2$, in verschiedenen Rollen gelesen.
Zähler: die nicht-trivialen $\mathbb{Z}_3$-Moden. Nenner: die Ordnung im
Quadrat, die Dimension des Operatorraums $\mathbb{C}^3\otimes\mathbb{C}^3$.

Das Entscheidende kommt danach. Das eine Matrixelement
$A = \langle v_0\mid R_5\,v_e\rangle$ hat Betrag $\sqrt2/3$ und Betragsquadrat
$2/9$. Die Koide-Formel
\[
  a_k = 1 + \sqrt2\cos\!\Bigl(\theta + \frac{2\pi k}{3}\Bigr)
\]
braucht genau diese zwei Zahlen: $\sqrt2$ als Koeffizient, $2/9$ als Phase. Mit
$3|A| = \sqrt2$ und $|A|^2 = 2/9$ wird daraus \BM:
\[
  a_k = 1 + 3|A|\cos\!\Bigl(|A|^2 + \frac{2\pi k}{3}\Bigr).
\]
Kein freier Parameter. Der Betrag gibt die Amplitude, das Betragsquadrat die
Phase -- Born-Regel-Struktur, hier aber nicht als Interpretation, sondern als
Algebra: dieselbe komplexe Zahl, einmal als $|A|$, einmal als $|A|^2$. Die
Koide-Formel ist damit vollständig aus der Frage abgeleitet: wie viel der
Fünffachdrehung landet in der symmetrischen $\mathbb{Z}_3$-Mode?

\section{Dok.~368: das $\phi$-Skelett}

Die beiden anderen Gewichte tragen die Fünf. Ihre Verhältnisse sind reine
$\phi$-Potenzen \BM:
\[
  \frac{p_1}{p_2} = \phi^8, \qquad \frac{p_0}{p_2} = 2\phi^4, \qquad
  3\sqrt{p_j} \in \{\sqrt2,\ \phi^2,\ \phi^{-2}\}.
\]
Und $\phi^8 \approx 46{,}98$ ist bis auf einen Galois-Faktor das Verhältnis
$m_\tau/m_e$: $74\,\phi^8$ trifft auf $0{,}03\,\%$, mit $\xi$-Korrektur auf
$0{,}003\,\%$ \KM. Das ist die zweite Rolle des Ikosaeders: $p_0$ liefert die
Phase, $p_1$ und $p_2$ liefern das Skelett der Massenverhältnisse. Beides aus
derselben Fünffachdrehung, beides ohne dass $\phi$ irgendwo hineingesteckt wurde.

\section{Dok.~293: warum es das Ikosaeder ist und nichts anderes}

Dok.~293 hat geprüft, ob $2/9$ an der Achse hängt. Zweihundert Zufallsachsen
mit Winkel $72^\circ$: keine liefert $2/9$. Winkelscan bei fester
Ikosaederachse: nur $72^\circ$ liefert $2/9$. Ordnungsreihe -- Drehungen der
Ordnung 2, 3, 4, 5, 6 um Ikosaederachsen: nur die Ordnung 5 liefert $2/9$ \KM.
Die Zahl ist ikosaeder-spezifisch. Sie entsteht, weil die Achse $(0,1,\phi)$
und der Winkel $72^\circ$ zusammen genau so zur $(1,1,1)$-Achse der
$\mathbb{Z}_3$ stehen, wie es die Ikosaedergeometrie vorschreibt.

\section{Was die ikosaedrische $\mathbb{Z}_3$ ist -- und was nicht}

Ein Vergleich schärft den Begriff. In der Stringtheorie gibt es Orbifolds
$T^6/\mathbb{Z}_3$, deren $\mathbb{Z}_3$ als $\omega\cdot\mathbb{1}$ auf
$\mathbb{C}^3$ wirkt -- alle drei Koordinaten bekommen dieselbe Phase. Diese
$\mathbb{Z}_3$ ist \emph{zentral}: sie vertauscht mit jeder Matrix. Die
ikosaedrische $\mathbb{Z}_3$ ist $C_3$ mit Eigenwerten $(1,\omega,\omega^2)$,
Spur null, \emph{nicht} zentral.

Beide Strukturen zählen $2/3^2$ -- die Orbifold-CFT über die Summe der
Twist-Feld-Gewichte $1/9+1/9$ \EM~(Dok.~367, externe Parallele). Aber nur die
ikosaedrische $\mathbb{Z}_3$ kann durch $R_5$ \emph{gemischt} werden. Bei der
zentralen wäre $\langle v_0|R_5|v_1\rangle$ null, weil $R_5$ dann mit
$\mathbb{Z}_3$ vertauscht und Sektoren erhält. Das Nicht-Vertauschen, das oben
die Zahl erzeugt hat, ist genau das, was die ikosaedrische Struktur von einem
Orbifold unterscheidet. $A_5$ hat triviales Zentrum; kein Element von $A_5$
kann in einer treuen Darstellung zu $\omega\cdot\mathbb{1}$ werden. Die beiden
$\mathbb{Z}_3$ sind verschiedene Objekte, die zufällig -- oder eben nicht
zufällig -- dieselbe Zählgröße haben.

\section{Was offen bleibt}

Zwei Dinge \SM.

Erstens die Galois-Faktoren. $\phi^8$ ist exakt, aber $m_\tau/m_e$ ist
$74\,\phi^8$, und $74 = 2\cdot37$ mit $37\mid|\text{GF}(3^{18})^*|$ kommt aus
der Zahlentheorie der Körper GF$(3^n)$ (Dok.~338), nicht aus der
Ikosaedergeometrie. Ob beide dieselbe Wurzel haben, ist nicht bekannt.

Zweitens die fraktale Korrektur. Der Schritt von Prozent- auf
$10^{-3}$-Niveau ist in allen Routen (Dok.~369) dieselbe Zahl in verschiedener
Kleidung -- gemessen, nicht hergeleitet. Die Koide-Route hat sie in
$\cos(2/9)$ versteckt und ist deshalb präzise; die algebraischen Routen müssen
sie nachtragen.

\section{Die kurze Fassung}

Das Ikosaeder ist der Ort, an dem Drei und Fünf sich nicht vertragen. FFGFT
bringt die Drei mit; das Ikosaeder ist die kleinste Struktur, die eine Fünf
hinzufügt, ohne dass die Fünf mit der Drei vertauscht. Aus diesem
Nicht-Vertauschen entsteht eine Mischung, und die Mischung hat eine rationale
Komponente -- $2/9$ -- und zwei $\phi$-Komponenten. Die rationale ist die
Koide-Phase. Die $\phi$-Komponenten sind das Skelett der Massen. Ein einziges
Matrixelement dieser Mischung trägt die gesamte Koide-Formel, parameterfrei.
